\scanpage{55}
% archive foliation: 54
Le cas 1°) \deleted{\unsure{dit}} signifie aussi que $T_\ell(E)(\beta)$ est tel que $\pi_1(X)$ y opère à travers un \uline{quotient} \uline{fini} [Mais cet énoncé ne se généralise guère via filtration de $\pi_1$, comme on voit sur $X=\mathrm{Spec}$ corps fini].

Supposons $E$ \inserted{$\in\mathrm{Ob}\,M^+$} \uline{semi-simple}, donc $T_\ell(E)$ semi-simple. Soit $E_i$ la composante de $E$ de type dimensionnel $i$, comment caractériser $T_\ell(E_i)\subset T_\ell(E)$ du pt de vue Galoisien ?
\inserted{Lorsque $X=\mathrm{Spec}$ corps fini, c'est purement \unsure{arithmétique} en termes de Frobenius.}

\deleted{
\unsure{les poids} \deleted{\unsure{homogènes de degré}} [la décomposition de $T_\ell(E)$ en composantes \deleted{\unsure{de degré}} homogènes est canoniquement déterminée]. On est ramené de proche en proche \deleted{[\unsure{Autant que}} \deleted{$\unsure{\mathbf{Z}(1)}$ et} [disons que le type dimensionnel précédent] à donner le maximum de type \inserted{dimensionnel} $i$, \unsure{par cas|quel cas} de type dimensionnel \deleted{\unsure{type dimensionnel 1}} \deleted{\unsure{alors}} \deleted{\unsure{le type dimensionnel}} \deleted{\unsure{qui sont des}} \unsure{qui sont nuls}.
}
\deleted{Le cas général \unsure{ne} \illegible{1} difficulté \unsure{pas} \illegible{2}, \unsure{si on connaît tous les Frobenius $\in\pi_1$}...}
